\chapter{Literature review}
\label{ch:literature}

The statistical literature connects four issues: information in financial news,
sentiment measurement, information lost through aggregation and the evidence
needed for prediction. A separate literature on risk controls explains why an
economic test must compare risk rules with controls that have similar market
exposure.

\section{Return information in news text}

News stories select facts and compare company results with expectations. A
sentiment model therefore measures framing, not fundamentals alone.

\citet{tetlock2007} finds that pessimism in a daily Wall Street Journal column
tracks downward market pressure and high volume before partially reversing. The
reversal distinguishes a temporary price move from a lasting change in value.
At the firm level, \citet{tetlock2008words} find that the fraction of negative
words predicts earnings and returns. Their measure counts words within a
document, whereas this dissertation counts negative stories within a firm-day.

The response to a price move depends on its source. \citet{chan2003news} finds
drift after identifiable public news, especially bad news, but reversal after
large returns without an attached story. \citet{savor2012information} reports a
similar split using analyst information. Negative coverage in this study could
therefore follow an earlier decline rather than precede a later one. The
recent-return controls in \cref{ch:methodology} test this reversal explanation.

Coverage and market conditions also change through time. \citet{fang2009} link
attention to expected returns, while \citet{engelberg2011} show that newspaper
distribution affects trading around earnings. \citet{garcia2013} finds that the
predictive content of news language is concentrated in recessions. Audience,
coverage and market state are therefore possible sources of temporal change.

Reuters studies motivate the separate LSEG analysis. Using 3.6 million articles,
\citet{uhl2014} finds greater index-level forecasting power in negative than
positive sentiment. With more than 900,000 Reuters stories,
\citet{heston2017} report daily predictability, weekly effects and slower
responses to negative news. Although their designs differ from this firm-day
comparison, both support asking whether an average sentiment score discards
useful information.

\section{Converting text to a numerical score}

The scoring method determines which language is positive, neutral or negative.
General-purpose dictionaries can misclassify financial terms that sound
negative in ordinary English but are neutral in company reporting.
\citet{loughran2011liability} discuss terms such as ``liability'' and ``tax'',
which a general dictionary can mark as negative even when they are routine in a
10-K. A finance-specific dictionary changes the labels and results.
\citet{loughran2016} also identify corpus design, negation, document length and
task-level validation as consequential choices.

Contextual models use the words around a phrase rather than counting words in
isolation. FinBERT adapts the BERT language model to financial text and classifies
sentences as positive, neutral or negative. ProsusAI FinBERT is trained further
on Financial PhraseBank, which records labels and disagreement among financially
trained readers \citep{malo2014,araci2019}. A separate finance-trained FinBERT
outperforms dictionaries and conventional machine learning in
\citet{huang2023}.

This study fixes one FinBERT version throughout. The model was trained on short
financial sentences, but it is applied here to FNSPID headlines, Reuters stories
and later periods outside that benchmark. Those differences create a measurement
limitation.

Prompted large language models can target firm-specific economic effects rather
than generic tone. \citet{kirtac2024} compare dictionaries, BERT-family models
and large language models on nearly a million US financial-news articles.
\citet{lopezlira2026} ask whether a headline is good or bad for the named firm
and find delayed responses concentrated in negative news and smaller firms,
with weaker results later. These studies motivate the prompt tests, which ask
whether a tailored score beats FinBERT on the same dates after costs.

Scoring old news with a modern model can mix look-ahead bias with distraction
from company knowledge. \citet{glasserman2023} compare original and
company-anonymised headlines to separate these effects and find that
anonymisation improves in-sample performance. Every prompted result here is
therefore a 2024 training-period experiment. Returns never enter the prompt,
and the experiment stops if none of the prompts meets the pre-specified
criteria.

\section{The information lost by averaging}

Each firm-day needs one signal, but document scores can be averaged, selected or
modelled separately. The studies reviewed here do not make a like-for-like
comparison of simple firm-day summaries. This study compares nine return-blind
measures using the same outcome, sample and adjustment for multiple tests.

Equal weighting gives a routine update the same weight as a first release and
allows positive or neutral coverage to offset a negative item. A threshold
measure keeps that negative item visible.

Earlier work also looks beyond average tone. \citet{allee2015} relate the
placement of positive and negative words within conference calls to firm
performance, reporting choices and user responses. \citet{isakin2023} link
news-sentiment dispersion to corporate-bond returns and uncertainty, while
\citet{chen2024} study dispersion around Chinese merger announcements. These
measures differ from a firm-day share but support retaining distributional
information.

Selection is another response to crowded news. \citet{uhl2021} construct a
selected ``top news'' measure and report that investors can benefit from
ignoring most news topics. This design instead keeps every story and fixes nine
summaries before a shared return test, so return relevance cannot determine
which stories enter.

Several stories may describe the same event. Stale and reprinted stories receive
a different market response and can be followed by reversal
\citep{tetlock2011stale}. Repeated updates about one adverse event also raise
negative share without adding information. Story count measures volume, but
FNSPID has no stable story-family or publisher identifier. The training
regression therefore cannot separate novelty from editorial repetition, which
limits both the association and its later non-replication.

\section{From a coefficient to an investment claim}

Training-period significance does not show that a result will persist. Many
proposed return predictors are unstable in later data \citep{welch2008}. Across
97 cross-sectional predictors, \citet{mclean2016} find smaller estimates outside
the original sample and after publication. Sampling error, selection and the
tendency to report unusually strong initial results can all contribute.

The order in which data become available also matters. In the exchange between
\citet{farmer2023} and \citet{cakici2025}, replacing a calculation that uses
future observations with one that uses only past observations sharply weakens
the reported forecasting value. A changed coefficient in this dissertation
shows non-replication, but cannot identify whether the market, sample or
measurement changed.

An information coefficient is the daily rank correlation between a signal and
returns. Information coefficients and portfolio returns answer different
questions: a precise association may still be too small to trade.
\citet{korajczyk2004} show that proportional costs and price impact materially
change momentum profitability and capacity, while \citet{novymarx2016} find the
worst damage among high-turnover anomalies. The economic tests report gross and
net return, turnover and per-side break-even cost.

When several predictors are tested, some may look significant by chance
\citep{harvey2016multiple}. The Benjamini--Hochberg (BH) adjustment limits the
expected share of false findings within a declared group of tests, under its
stated assumptions \citep{benjamini1995}. It covers the nine-rule comparison and
related groups, not specifications designed after seeing results. The training
coefficient is therefore treated as exploratory.

Nearby sessions can be related, and their variance can change. Newey--West
heteroskedasticity and autocorrelation consistent (HAC) standard errors allow
for both features
\citep{newey1987}. Portfolio comparisons use fixed-length block resampling to
keep nearby sessions together \citep{kunsch1989}; this implementation wraps the
blocks at the sample boundary. It differs from the random-length stationary
bootstrap of \citet{politis1994}.

An interval crossing zero can reflect either a small effect or low precision.
Minimum detectable effects (MDEs) quantify the effect size a test has a good
chance of detecting \citep{bloom1995}. If a short LSEG sample can detect only
coefficients several times larger than the FNSPID estimate, its null cannot
establish whether an FNSPID-sized relationship exists in LSEG\@.

\section{Evaluating risk overlays}

A weak directional signal might still identify periods for lower risk. However,
time in cash mechanically lowers volatility and losses, so the control must have
similar market exposure.

The sentiment-free benchmark is a heterogeneous autoregressive (HAR) volatility
forecast. \citet{corsi2009har} models realised volatility with daily, weekly and
monthly components. In the benchmark used here, exposure falls as the forecast
rises. Evidence on inverse-volatility factor scaling \citep{moreira2017}
motivates this benchmark without endorsing sentiment timing.

The first comparator is the fixed HAR path. The second holds a constant
fraction of that path to match the news rule's average exposure. In FNSPID, the
constant is calibrated after observing the 2020--2023 testing results, so it supports only
descriptive lower-exposure attribution. In LSEG, matched constants are selected
in the training folds before forward testing.

A circular-shift check moves the full risk-off schedule through returns while
preserving its activation frequency. It tests whether the same-session risk
schedule aligns unusually with downside outcomes, but not whether news was
observable before returns began. A timing claim needs both unusual alignment
and verified pre-return availability. The matched constant separately tests
whether lower average exposure explains the result.

\section{The research gap}

The measurement comparison uses one outcome, a declared test family and
dependence-aware uncertainty. An investment claim also needs testing-period
evidence and trading costs. This dissertation applies those requirements to a
hard-threshold firm-day measure through training selection, timing, testing,
economic and separate-source checks. The analysis retains null results without
using them to start a further search.
